\documentclass[12pt,a4paper]{article}

\usepackage[utf8]{inputenc}
\usepackage[T1]{fontenc}
\usepackage[brazil]{babel}
\usepackage[margin=2.5cm]{geometry}
\usepackage{amsmath,amssymb,amsthm}
\usepackage{booktabs}
\usepackage{graphicx}
\usepackage{xcolor}
\usepackage{tikz}
\usepackage[hidelinks]{hyperref}

\newtheorem{teorema}{Teorema}
\newtheorem{proposicao}{Proposição}
\theoremstyle{definition}
\newtheorem{definicao}{Definição}

% Marcação de partes ainda não escritas
\newcommand{\afazer}[1]{\textcolor{red}{\textbf{[A FAZER: #1]}}}

\title{Sistema de Recomendação Baseado em Grafos Bipartidos}
\author{Alice Barros Lourencini Palaoro \and Felipe Augusto da Silva Cruz \and
        Gabriel Andrade Souza \and Guilherme Henrique Pereira Cunha}
\date{Matemática Discreta --- Novembro de 2026}

\begin{document}

\maketitle

\begin{abstract}
\afazer{resumo de 5--8 linhas, escrever por último: problema, modelo (grafo
bipartido e contagem de passeios), base MovieLens, principais resultados.}
\end{abstract}

% =====================================================================
\section{Introdução}
% =====================================================================


\indent Atualmente, os algoritmos de recomendação desempenham um papel crucial em praticamente qualquer tipo de produto ou serviço digital. Tendo isso em vista, o presente projeto propõe a construção de uma solução baseada em grafos, em que os vértices representam usuários, produtos e categorias, e as arestas indicam as interações entre eles (como usuário-produto e produto-categoria). Para validar a abordagem, a implementação utilizou como estudo de caso a indústria de cinema e streaming, empregando uma base de dados de avaliações de usuários com o objetivo de recomendar títulos inéditos por meio de passeios estruturados no grafo.
A ideia central do trabalho consiste em mapear e analisar os caminhos formados no grafo entre o usuário e os itens ainda não avaliados por ele. A partir das propriedades e características desses trajetos, o algoritmo pontua e seleciona os filmes mais adequados, baseando-se no histórico de consumo e nas preferências anteriores do indivíduo.
Por fim, vale destacar que, apesar do domínio escolhido para este estudo, a arquitetura desenvolvida mantém sua relevância e versatilidade, podendo ser aplicada a qualquer ecossistema estruturado em usuários, itens e avaliações. Exemplos disso incluem plataformas de e-commerce de vestuário, serviços de streaming de música, livrarias digitais e marketplaces de eletrônicos, o que evidencia o alto potencial de generalização e escalabilidade do modelo proposto.


% =====================================================================
\section{Modelagem}
% =====================================================================

\subsection{Grafo bipartido de interações}

\begin{definicao}
Um grafo $G = (V, E)$ é \emph{bipartido} se existe uma partição
$V = U \cup I$, com $U \cap I = \emptyset$, tal que toda aresta de $E$ tem
uma extremidade em $U$ e a outra em $I$.
\end{definicao}

Modelamos o problema com $U$ sendo o conjunto de usuários, $I$ o conjunto de
itens (filmes) e uma aresta $\{u, i\} \in E$ sempre que o usuário $u$ avaliou
o filme $i$. Não há arestas entre dois usuários nem entre dois filmes.

\subsection{Matriz de biadjacência}

\begin{definicao}
Fixadas bijeções $U \to \{0, \dots, |U|-1\}$ e $I \to \{0, \dots, |I|-1\}$, a
\emph{matriz de biadjacência} de $G$ é a matriz $B \in \{0,1\}^{|U| \times |I|}$
dada por
\[
  B_{ui} =
  \begin{cases}
    1, & \text{se } \{u, i\} \in E,\\
    0, & \text{caso contrário.}
  \end{cases}
\]
\end{definicao}

Diferente da matriz de adjacência usual, que é quadrada de ordem $|U|+|I|$,
$B$ é retangular. As duas se relacionam por
\[
  A = \begin{pmatrix} 0 & B \\ B^{\mathsf T} & 0 \end{pmatrix}.
\]

A soma da linha $u$ de $B$ é o grau de $u$ (quantos filmes ele avaliou) e a
soma da coluna $i$ é o grau de $i$ (quantos usuários o avaliaram). Como cada
aresta tem exatamente uma extremidade em cada lado, vale a versão bipartida
do lema do aperto de mãos:
\[
  \sum_{u \in U} \deg(u) = \sum_{i \in I} \deg(i) = |E|.
\]

\subsection{Produtos de $B$ e contagem de passeios}

\paragraph{$BB^{\mathsf T}$ (usuário $\times$ usuário).}
\[
  (BB^{\mathsf T})_{uv} = \sum_{i \in I} B_{ui} B_{vi}.
\]
O produto $B_{ui}B_{vi}$ vale 1 exatamente quando $u$ e $v$ avaliaram $i$.
Logo $(BB^{\mathsf T})_{uv}$ é o número de filmes em comum entre $u$ e $v$,
que é também o número de passeios de comprimento 2 de $u$ a $v$
($u \to i \to v$). A diagonal é o grau do usuário, e a matriz é simétrica,
pois $(BB^{\mathsf T})^{\mathsf T} = BB^{\mathsf T}$.

\paragraph{$BB^{\mathsf T}B$ (usuário $\times$ filme).}
\[
  S_{ui} = (BB^{\mathsf T}B)_{ui} = \sum_{v \in U} (BB^{\mathsf T})_{uv}\, B_{vi}.
\]
$S_{ui}$ conta os passeios de comprimento 3 da forma
$u \to i' \to v \to i$: ``você viu $i'$, o usuário $v$ também viu $i'$, e $v$
viu $i$''. Cada usuário $v$ que viu $i$ contribui com peso igual ao número de
filmes que tem em comum com $u$. É uma forma de filtragem colaborativa obtida
apenas por contagem.

\begin{teorema}\label{teo:passeios}
Seja $A$ a matriz de adjacência de um grafo. Para todo $k \geq 1$,
$(A^k)_{xy}$ é o número de passeios de comprimento $k$ de $x$ a $y$.
\end{teorema}

\begin{proof}
\afazer{prova por indução em $k$. Base: $k = 1$ (definição de $A$). Passo:
$(A^{k+1})_{xy} = \sum_z (A^k)_{xz} A_{zy}$; todo passeio de comprimento
$k+1$ de $x$ a $y$ é um passeio de comprimento $k$ de $x$ a algum $z$
seguido da aresta $\{z, y\}$. Depois, mostrar que $(A^3)$ restrita ao bloco
$U \times I$ é exatamente $BB^{\mathsf T}B$.}
\end{proof}

\paragraph{Passeios e caminhos.}
O teorema conta \emph{passeios}, que podem repetir vértices. Por exemplo,
se $u$ viu $i$, o passeio $u \to i \to u \to i$ existe; logo $S_{ui} \geq 1$
para todo filme já visto. Por isso esses filmes são excluídos da
recomendação. Já para um filme $i$ que $u$ \emph{não} viu, todo passeio
$u \to i' \to v \to i$ é um caminho simples: $v \neq u$ (pois $v$ viu $i$ e
$u$ não) e $i' \neq i$ (pois $u$ viu $i'$ e não viu $i$). Assim, nos itens
candidatos à recomendação, passeios e caminhos coincidem.

\paragraph{Regra de recomendação.}
Para cada usuário $u$, recomendamos os $k$ filmes $i$ com $B_{ui} = 0$ de maior
$S_{ui}$. A pontuação não é uma nota prevista nem uma probabilidade; serve
apenas para ordenar.

\subsection{Grafo tripartido: incluindo as tags}

O modelo bipartido só enxerga quem assistiu o quê. Para usar também o
\emph{conteúdo} dos filmes, acrescentamos um terceiro conjunto de vértices
$T$, formado pelas tags (etiquetas como ``sci-fi'', ``romance'',
``pixar''), com uma aresta $\{i, t\}$ sempre que o filme $i$ recebeu a
tag $t$. O grafo passa a ter vértices $U \cup I \cup T$, e as arestas só
ligam usuário--filme ou filme--tag.

\begin{definicao}
A \emph{matriz filme--tag} é $M \in \{0,1\}^{|I| \times |T|}$, com
\[
  M_{it} =
  \begin{cases}
    1, & \text{se o filme } i \text{ recebeu a tag } t,\\
    0, & \text{caso contrário.}
  \end{cases}
\]
\end{definicao}

% Obs.: chamamos de M (e não de G ou T) porque G já é o nome do grafo
% e T é o nome do conjunto de tags. Troquem se o grupo preferir outra letra.

\paragraph{$MM^{\mathsf T}$ (filme $\times$ filme).}
\[
  (MM^{\mathsf T})_{ij} = \sum_{t \in T} M_{it} M_{jt}.
\]
O termo $M_{it}M_{jt}$ vale 1 quando os filmes $i$ e $j$ têm a tag $t$.
Assim, $(MM^{\mathsf T})_{ij}$ é o número de tags em comum entre $i$ e $j$,
ou seja, o número de passeios $i \to t \to j$ de comprimento 2. É uma
medida de semelhança de \emph{conteúdo} entre filmes.

\paragraph{$BMM^{\mathsf T}$ (usuário $\times$ filme).}
\[
  (BMM^{\mathsf T})_{ui} = \sum_{i' \in I} B_{ui'}\,(MM^{\mathsf T})_{i'i}.
\]
Cada filme $i'$ que $u$ viu contribui com o número de tags que ele
compartilha com $i$. Logo $(BMM^{\mathsf T})_{ui}$ conta os passeios de
comprimento 3 da forma $u \to i' \to t \to i$: ``você viu $i'$, que tem a
tag $t$, e $i$ também tem a tag $t$''. Enquanto $BB^{\mathsf T}B$ é
filtragem \emph{colaborativa} (passa por outros usuários),
$BMM^{\mathsf T}$ é filtragem \emph{por conteúdo} (passa pelas tags).

\paragraph{Pontuação híbrida.}
As duas matrizes têm o mesmo formato $|U| \times |I|$, então podem ser
somadas:
\[
  S = BB^{\mathsf T}B + \alpha\, BMM^{\mathsf T}, \qquad \alpha \geq 0.
\]
$S_{ui}$ soma os passeios $u \to i' \to v \to i$ (via usuários) com os
passeios $u \to i' \to t \to i$ (via tags), estes últimos com peso
$\alpha$. O parâmetro controla o equilíbrio: $\alpha = 0$ recupera o
modelo bipartido, e $\alpha$ grande faz o conteúdo dominar. Como as duas
parcelas contam coisas diferentes e podem ter ordens de grandeza
diferentes, $\alpha$ também serve para colocá-las na mesma escala. A regra
de recomendação é a mesma: os $k$ filmes não vistos de maior $S_{ui}$.

\paragraph{Exemplo.}
No exemplo da Seção~3, considere as tags \emph{ação}, \emph{romance} e
\emph{animação}: Matrix tem \emph{ação}; Titanic, \emph{romance};
Toy Story, \emph{animação}; e Up, \emph{romance} e \emph{animação}.
O grafo tripartido resultante está na Figura~\ref{fig:tripartido}.

\begin{figure}[ht]
\centering
\begin{tikzpicture}[
    usr/.style={circle, draw, fill=blue!15, minimum size=1.1cm, font=\small},
    itm/.style={rectangle, draw, fill=orange!20, minimum width=2cm,
                minimum height=0.7cm, font=\small},
    tag/.style={rectangle, rounded corners, draw, fill=green!15,
                minimum width=1.8cm, minimum height=0.7cm, font=\small}]
  \node[usr] (ana)   at (0, 0)    {Ana};
  \node[usr] (bruno) at (0, -1.6) {Bruno};
  \node[usr] (carla) at (0, -3.2) {Carla};
  \node[itm] (mat) at (4.5, 0.4)   {Matrix};
  \node[itm] (tit) at (4.5, -0.9)  {Titanic};
  \node[itm] (toy) at (4.5, -2.2)  {Toy Story};
  \node[itm] (up)  at (4.5, -3.5)  {Up};
  \node[tag] (aca) at (9, 0)    {ação};
  \node[tag] (rom) at (9, -1.6) {romance};
  \node[tag] (ani) at (9, -3.2) {animação};
  \draw (ana) -- (mat);   \draw (ana) -- (tit);
  \draw (bruno) -- (mat); \draw (bruno) -- (toy); \draw (bruno) -- (up);
  \draw (carla) -- (tit); \draw (carla) -- (toy);
  \draw[dashed] (mat) -- (aca); \draw[dashed] (tit) -- (rom);
  \draw[dashed] (toy) -- (ani); \draw[dashed] (up) -- (rom);
  \draw[dashed] (up) -- (ani);
\end{tikzpicture}
\caption{Grafo tripartido do exemplo: usuários, filmes e tags. As arestas
tracejadas são as da matriz $M$.}
\label{fig:tripartido}
\end{figure}

Com linhas (Matrix, Titanic, Toy Story, Up) e colunas
(ação, romance, animação):
\[
  M = \begin{pmatrix}
    1 & 0 & 0\\
    0 & 1 & 0\\
    0 & 0 & 1\\
    0 & 1 & 1
  \end{pmatrix},
  \qquad
  MM^{\mathsf T} = \begin{pmatrix}
    1 & 0 & 0 & 0\\
    0 & 1 & 0 & 1\\
    0 & 0 & 1 & 1\\
    0 & 1 & 1 & 2
  \end{pmatrix},
  \qquad
  BMM^{\mathsf T} = \begin{pmatrix}
    1 & 1 & \mathbf{0} & \mathbf{1}\\
    1 & \mathbf{1} & 2 & 3\\
    \mathbf{0} & 1 & 1 & \mathbf{2}
  \end{pmatrix}.
\]
Para Carla, o modelo bipartido dava \mbox{Matrix $= 2$} e \mbox{Up $= 1$}. A parcela de
tags dá \mbox{Matrix $= 0$} e \mbox{Up $= 2$}, pois Up compartilha \emph{romance} com
Titanic e \emph{animação} com Toy Story, filmes que Carla viu. Com
$\alpha = 1$, a pontuação final fica \mbox{Matrix $= 2$} e \mbox{Up $= 3$}, e a
recomendação de Carla muda para \emph{Up}: um filme parecido com o que ela
gosta, e não apenas um filme visto por usuários parecidos com ela.

% =====================================================================
\section{Exemplo ilustrativo}
% =====================================================================

Considere três usuários e quatro filmes: Ana viu \emph{Matrix} e
\emph{Titanic}; Bruno viu \emph{Matrix}, \emph{Toy Story} e \emph{Up}; Carla
viu \emph{Titanic} e \emph{Toy Story} (Figura~\ref{fig:exemplo}).

\begin{figure}[ht]
\centering
\begin{tikzpicture}[
    usr/.style={circle, draw, fill=blue!15, minimum size=1.1cm, font=\small},
    itm/.style={rectangle, draw, fill=orange!20, minimum width=2cm,
                minimum height=0.7cm, font=\small}]
  \node[usr] (ana)   at (0, 0)    {Ana};
  \node[usr] (bruno) at (0, -1.6) {Bruno};
  \node[usr] (carla) at (0, -3.2) {Carla};
  \node[itm] (mat) at (5, 0.4)   {Matrix};
  \node[itm] (tit) at (5, -0.9)  {Titanic};
  \node[itm] (toy) at (5, -2.2)  {Toy Story};
  \node[itm] (up)  at (5, -3.5)  {Up};
  \draw (ana) -- (mat);   \draw (ana) -- (tit);
  \draw (bruno) -- (mat); \draw (bruno) -- (toy); \draw (bruno) -- (up);
  \draw (carla) -- (tit); \draw (carla) -- (toy);
\end{tikzpicture}
\caption{Grafo bipartido do exemplo: usuários à esquerda, filmes à direita.}
\label{fig:exemplo}
\end{figure}

Com as linhas na ordem (Ana, Bruno, Carla) e as colunas na ordem
(Matrix, Titanic, Toy Story, Up):
\[
  B = \begin{pmatrix}
    1 & 1 & 0 & 0\\
    1 & 0 & 1 & 1\\
    0 & 1 & 1 & 0
  \end{pmatrix},
  \qquad
  BB^{\mathsf T} = \begin{pmatrix}
    2 & 1 & 1\\
    1 & 3 & 1\\
    1 & 1 & 2
  \end{pmatrix},
  \qquad
  BB^{\mathsf T}B = \begin{pmatrix}
    3 & 3 & \mathbf{2} & \mathbf{1}\\
    4 & \mathbf{2} & 4 & 3\\
    \mathbf{2} & 3 & 3 & \mathbf{1}
  \end{pmatrix}.
\]
Em negrito estão os filmes ainda não vistos. Por exemplo,
$S_{\text{Ana},\,\text{Toy Story}} = 2$ corresponde aos caminhos
Ana $\to$ Matrix $\to$ Bruno $\to$ Toy Story e
Ana $\to$ Titanic $\to$ Carla $\to$ Toy Story. As recomendações são:
Ana $\to$ \emph{Toy Story}, Bruno $\to$ \emph{Titanic} e
Carla $\to$ \emph{Matrix}. Os valores foram calculados à mão e conferidos
com o script \texttt{src/exemplo\_pequeno.py}.

% =====================================================================
\section{Base de dados}
% =====================================================================

Usamos a base \emph{MovieLens latest-small}, disponibilizada pelo GroupLens, com avaliações de 0{,}5 a 5{,}0 estrelas feitas
por usuários do site MovieLens. Usamos o arquivo \texttt{ratings.csv} (pares
usuário--filme--nota) e o \texttt{movies.csv} (títulos). As estatísticas da
Tabela~\ref{tab:base} foram obtidas pelo script \texttt{src/explorar.py}.

\begin{table}[ht]
\centering
\begin{tabular}{lr}
\toprule
Medida & Valor \\
\midrule
Usuários $|U|$ & 610 \\
Filmes $|I|$ (com ao menos uma avaliação) & 9\,724 \\
Arestas $|E|$ & 100\,836 \\
Pares (usuário, filme) repetidos & 0 \\
Grau dos usuários: mín.\ / mediana / média / máx. & 20 / 70{,}5 / 165{,}3 / 2\,698 \\
Grau dos filmes: mín.\ / mediana / média / máx. & 1 / 3 / 10{,}4 / 329 \\
Filmes com grau 1 & 3\,446 (35\%) \\
Densidade $|E| / (|U|\,|I|)$ & 1{,}7\% \\
\bottomrule
\end{tabular}
\caption{Estatísticas do grafo bipartido construído a partir do MovieLens.}
\label{tab:base}
\end{table}

Observações:
\begin{itemize}
  \item Não há pares repetidos, então cada linha do arquivo é uma aresta
        distinta e $B$ tem apenas entradas 0 e 1. O lema do aperto de mãos
        foi verificado: as duas somas de graus valem $100\,836 = |E|$.
  \item O \texttt{movies.csv} lista 9\,742 filmes, mas 18 nunca foram
        avaliados; por serem vértices isolados, ficam fora do grafo.
  \item Todo usuário tem grau pelo menos 20, pois a base inclui apenas
        usuários com no mínimo 20 avaliações.
  \item Os graus dos filmes têm cauda longa: a média é cerca de 10, mas a
        mediana é 3. Poucos filmes concentram muitas avaliações
        (\emph{Forrest Gump}: 329; \emph{Um Sonho de Liberdade}: 317;
        \emph{Pulp Fiction}: 307).
  \item A densidade de 1{,}7\% significa que 98{,}3\% das
        $610 \times 9\,724 \approx 5{,}9$ milhões de entradas de $B$ são
        zero, o que justifica o armazenamento esparso.
\end{itemize}

% =====================================================================
\section{Implementação}
% =====================================================================

O código está em Python, usando \texttt{pandas}, \texttt{numpy} e
\texttt{scipy.sparse}.

\paragraph{Construção de $B$ (\texttt{src/matriz.py}).}
Os identificadores da base não são contíguos (\texttt{movieId} vai até
193\,609, com apenas 9\,724 filmes). Por isso construímos explicitamente as
bijeções da Definição de $B$: os identificadores distintos, em ordem
crescente, recebem os índices $0, 1, \dots$. Cada aresta $(u, i)$ é então
traduzida para a posição (linha, coluna), e $B$ é montada no formato esparso
CSR (\emph{Compressed Sparse Row}), que armazena apenas as $|E|$ entradas não
nulas. Conferimos que $B$ tem formato $610 \times 9\,724$, exatamente
$100\,836$ entradas não nulas e que as somas de linhas e colunas reproduzem
os graus da Tabela~\ref{tab:base}.

\paragraph{Ordem das multiplicações.}
Pela associatividade, $(BB^{\mathsf T})B = B(B^{\mathsf T}B)$, mas o custo é
diferente: $BB^{\mathsf T}$ tem formato $610 \times 610$ (cerca de 372 mil
entradas), enquanto $B^{\mathsf T}B$ tem formato $9\,724 \times 9\,724$
(cerca de 94{,}5 milhões). Calculamos da esquerda para a direita.

\paragraph{Esparsidade de $S$.}
Embora $B$ seja esparsa, $S = BB^{\mathsf T}B$ tem 99{,}1\% das entradas não
nulas: em três passos, quase todo usuário alcança quase todo filme. Por isso
$S$ é convertida para uma matriz densa antes da ordenação.

\paragraph{Recomendação (\texttt{src/recomendar.py}).}
As posições dos filmes já vistos (as entradas não nulas de $B$) são zeradas
em $S$, e para cada linha são tomados os 10 índices de maior pontuação. O
cálculo completo leva menos de um segundo.

% =====================================================================
\section{Resultados preliminares}
% =====================================================================

A Tabela~\ref{tab:top10} mostra as recomendações para o usuário 1.

\begin{table}[ht]
\centering
\begin{tabular}{rlr}
\toprule
\# & Filme & $S_{ui}$ \\
\midrule
1  & The Shawshank Redemption (1994) & 10\,745 \\
2  & Terminator 2: Judgment Day (1991) & 10\,205 \\
3  & The Sixth Sense (1999) & 9\,323 \\
4  & The Godfather (1972) & 8\,942 \\
5  & Die Hard (1988) & 8\,541 \\
6  & The Lord of the Rings: The Fellowship of the Ring (2001) & 8\,461 \\
7  & Speed (1994) & 8\,421 \\
8  & Twelve Monkeys (1995) & 8\,402 \\
9  & True Lies (1994) & 8\,225 \\
10 & Aladdin (1992) & 8\,044 \\
\bottomrule
\end{tabular}
\caption{Top-10 por $BB^{\mathsf T}B$ para o usuário 1.}
\label{tab:top10}
\end{table}

As recomendações são plausíveis, mas pouco personalizadas
(Tabela~\ref{tab:vies}).

\begin{table}[ht]
\centering
\begin{tabular}{lr}
\toprule
Medida & Valor \\
\midrule
Filmes distintos somando os top-10 dos 610 usuários & 153 \\
Sobreposição média com o top-10 por popularidade & 6{,}6 de 10 \\
Grau mediano dos filmes recomendados & 211 \\
Correlação entre $S_{ui}$ e $\deg(i)$ (usuário 1) & 0{,}96 \\
Usuários que recebem \emph{Star Wars: Episode IV} & 331 \\
\bottomrule
\end{tabular}
\caption{Indicadores de viés de popularidade do recomendador $BB^{\mathsf T}B$.}
\label{tab:vies}
\end{table}

A explicação vem da própria fórmula: em
$S_{ui} = \sum_v (BB^{\mathsf T})_{uv} B_{vi}$, um filme com muitos
espectadores tem muitos termos não nulos, e os usuários de maior peso
(muitos filmes em comum com $u$) tendem a ser os que viram muitos filmes,
inclusive os populares. Vértices de grau alto atraem passeios. Esse viés
motiva a comparação com um \emph{baseline} de popularidade e a normalização
pelo grau, descritas a seguir.

% =====================================================================
\section{Baseline de popularidade}
% =====================================================================

\afazer{etapa 7: recomendar os filmes de maior grau não vistos.}

% =====================================================================
\section{Avaliação}
% =====================================================================

\paragraph{Protocolo.}
Para cada usuário, separamos aleatoriamente 20\% das avaliações para
\emph{teste} e usamos o restante (\emph{treino}) para montar as matrizes
$B$ e $M$. Cada modelo gera, só com o treino, uma lista $R_u$ de $K = 10$
filmes para cada usuário $u$. Consideramos \emph{relevantes} os filmes do
teste com nota pelo menos $4{,}0$:
\[
  \mathrm{Rel}_u = \{\, i : (u, i) \text{ está no teste e a nota é } \geq 4{,}0 \,\}.
\]
As médias são tomadas sobre o conjunto $U'$ dos usuários com
$\mathrm{Rel}_u \neq \emptyset$, e $\deg(i)$ denota o grau do filme $i$ no
grafo de treino.

\subsection{Métricas de acerto}

\paragraph{Precision@$K$.} Fração da lista que acertou:
\[
  \mathrm{Precision@}K = \frac{1}{|U'|} \sum_{u \in U'}
  \frac{|R_u \cap \mathrm{Rel}_u|}{K}.
\]

\paragraph{Recall@$K$.} Fração dos relevantes que apareceram na lista:
\[
  \mathrm{Recall@}K = \frac{1}{|U'|} \sum_{u \in U'}
  \frac{|R_u \cap \mathrm{Rel}_u|}{|\mathrm{Rel}_u|}.
\]

\paragraph{Hit rate@$K$.} Fração dos usuários com pelo menos um acerto:
\[
  \mathrm{HR@}K = \frac{1}{|U'|} \sum_{u \in U'}
  \mathbf{1}\bigl[\, R_u \cap \mathrm{Rel}_u \neq \emptyset \,\bigr].
\]

\subsection{Métricas de viés de popularidade}

As três métricas acima medem apenas se o modelo acerta. Como vimos na
Seção de resultados preliminares, recomendar sempre os filmes mais
populares já garante muitos acertos. Por isso usamos também três métricas
que medem \emph{o que} está sendo recomendado.

\paragraph{Recall de baixo grau.}
Seja $I_{\text{baixo}}$ o conjunto dos filmes de \emph{cauda longa}: os que
ficam fora dos 20\% de maior grau no treino. O recall de baixo grau é o
recall calculado só sobre esses filmes:
\[
  \mathrm{Recall}_{\text{baixo}}\mathrm{@}K = \frac{1}{|U''|} \sum_{u \in U''}
  \frac{|R_u \cap \mathrm{Rel}_u \cap I_{\text{baixo}}|}
       {|\mathrm{Rel}_u \cap I_{\text{baixo}}|},
\]
onde $U''$ são os usuários com $\mathrm{Rel}_u \cap I_{\text{baixo}} \neq
\emptyset$. Ele mede se o modelo consegue acertar filmes pouco conhecidos,
que são justamente os que o usuário teria mais dificuldade de encontrar
sozinho. Um recomendador por popularidade tem recall de baixo grau
praticamente nulo.

\paragraph{Cobertura.}
Fração do catálogo que aparece em alguma lista:
\[
  \mathrm{Cobertura@}K = \frac{\bigl|\bigcup_{u \in U} R_u\bigr|}{|I|}.
\]
Em termos do grafo, é a fração de vértices de $I$ alcançados pelas
recomendações. Com o $BB^{\mathsf T}B$ binário, os top-10 dos 610 usuários
somam apenas 153 filmes distintos (Tabela~\ref{tab:vies}), uma cobertura
de $153/9\,724 \approx 1{,}6\%$.

\paragraph{Grau médio.}
Grau médio dos filmes recomendados:
\[
  \overline{\deg}\mathrm{@}K = \frac{1}{|U|\,K} \sum_{u \in U} \sum_{i \in R_u} \deg(i).
\]
Quanto maior, mais o modelo se concentra em vértices de grau alto, ou
seja, em filmes populares. Serve para medir diretamente o efeito descrito
nos resultados preliminares: vértices de grau alto atraem passeios.

\paragraph{Leitura conjunta.}
Para as métricas de acerto e para o recall de baixo grau e a cobertura,
maior é melhor. O grau médio não tem um valor ``certo'', mas um valor
menor com acerto parecido indica recomendações mais personalizadas. Em
geral há um compromisso: reduzir o viés de popularidade costuma custar um
pouco de precision, e as seis métricas juntas mostram esse equilíbrio.

\subsection{Resultados}

\begin{table}[ht]
\centering
\begin{tabular}{lcccccc}
\toprule
Modelo & Prec@10 & Rec@10 & HR@10 & Rec$_{\text{baixo}}$@10 & Cob@10 & $\overline{\deg}$@10 \\
\midrule
Popularidade          & -- & -- & -- & -- & -- & -- \\
Bipartido binário     & -- & -- & -- & -- & -- & -- \\
Bipartido ponderado   & -- & -- & -- & -- & -- & -- \\
Tripartido ponderado  & -- & -- & -- & -- & -- & -- \\
\bottomrule
\end{tabular}
\caption{Comparação dos modelos ($K = 10$, $\alpha = 0{,}5$).}
\label{tab:avaliacao}
\end{table}

\afazer{preencher a Tabela~\ref{tab:avaliacao} com a saída de
\texttt{src/benchmark.py} e comentar os resultados.}

% =====================================================================
\section{Melhoria: normalização pelo grau}
% =====================================================================

\afazer{etapa 9: passeio aleatório / normalização pelo grau
ou PageRank personalizado; repetir a avaliação e a
Tabela~\ref{tab:vies}.}

% =====================================================================
\section{Conclusão}
% =====================================================================

\afazer{conclusão.}


\end{document}
